% TOP_PROBLEM: {"problemId": "problem.top500-omission-w045049-f-conjecture-mbar-zero-n", "canonicalTitle": "F-Conjecture for the Nef Cone of Mbar_(0,n)", "releaseRank": 499, "primaryDomain": "geometry_topology", "primaryDomainLabel": "Geometry & topology", "displayStatus": "open_with_solved_subcases", "releaseStatus": "open_with_solved_subcases"}
% !TeX program = lualatex
% Definition and sources only; this file is not an LLM solution attempt.
\documentclass[11pt]{article}
\usepackage[margin=25mm]{geometry}
\usepackage{fontspec}
\setmainfont{Latin Modern Roman}
\usepackage{amsmath,amssymb,mathtools}
\usepackage{unicode-math}
\setmathfont{Latin Modern Math}
\usepackage{xurl,hyperref}
\hypersetup{colorlinks=true,urlcolor=blue}
\setcounter{secnumdepth}{0}
\setlength{\parindent}{0pt}
\setlength{\parskip}{5pt}
\setlength{\emergencystretch}{3em}
\begin{document}
\section*{\texorpdfstring{F-Conjecture for the Nef Cone of Mbar\_(0,n)}{F-Conjecture for the Nef Cone of Mbar\_(0,n)}}\label{499}
\subsection{Definitions and mathematical statement}
The projective moduli space $\overline M_{0,n}$ parametrizes stable genus-zero nodal curves with $n$ distinct labeled smooth marked points. Stability requires at least three marked points or nodes on each irreducible component. An F-curve is the one-dimensional family obtained by partitioning the labels into four nonempty groups, attaching the corresponding fixed tails, and varying the cross-ratio on the central four-pointed component, including its stable limits.

For a real numerical divisor class $D$, nefness means $D\cdot C\ge0$ for every complete curve $C$. The conjecture says
\[
 D\text{ nef}\quad\Longleftrightarrow\quad
 D\cdot F\ge0\text{ for every F-curve }F,
\]
or dually that $\overline{\operatorname{NE}}_1(\overline M_{0,n})$ is generated by their classes. This is for every $n\ge3$. For $n=3$ the space is a point; for $n=4$ the relevant F-curve is $\overline M_{0,4}$ itself. Thus the phrase boundary stratum uses the conventional extension to this case, not a requirement that the curve lie in the proper boundary.
\par\smallskip\textit{Formulation and concept references:} \href{https://arxiv.org/html/2507.12434v1}{[S1]}.

\subsection{Short English statement}
Can nef divisors on every genus-zero pointed-curve moduli space be recognized by testing only the standard one-dimensional F-curves?
\subsection{Sources}
\begin{itemize}
\item[S1]M. Fedorchuk and A. Mellit, Symmetric and non-symmetric F-conjectures are equivalent, arXiv 2507.12434.
\url{https://arxiv.org/html/2507.12434v1}.
\textit{Formulation source.}
\end{itemize}
\end{document}
